\documentclass[12pt]{article}
\usepackage[T1]{fontenc}
\usepackage[utf8]{inputenc}
\usepackage{lmodern}
\usepackage{textcomp}
\usepackage{enumitem}
\usepackage{geometry}
\usepackage{graphicx}
\usepackage{amsmath, amssymb}
\geometry{margin=1in}
\begin{document}
\begin{center}
Sociology - Graduate Level Assessment 7 \
Total Marks: 40 \
Answer all questions.
\end{center}

\section*{Multiple Choice (10 marks)}
\begin{enumerate}[label=\arabic*.]
\item Urbanization occurred in the nineteenth century because:
\begin{enumerate}[label=(\alph*)]
    \item commuters started moving out of villages and into cities
    \item towns and cities were becoming increasingly planned and managed
    \item industrial capitalism led to a shift of population from rural to urban areas
    \item transport systems were not provided, so it was easier to live in the city
\end{enumerate}
\item Mass-society theory suggests that:
\begin{enumerate}[label=(\alph*)]
    \item the content of the media is determined by market forces
    \item the subordinate classes are dominated by the ideology of the ruling class
    \item the media manipulate 'the masses' as vulnerable, passive consumers
    \item audiences make selective interpretations of media messages
\end{enumerate}
\item The 'two-sex' model that Laqueur (1990) identified:
\begin{enumerate}[label=(\alph*)]
    \item contrasted homosexuality with heterosexuality
    \item distinguished between male and females as separate sexes
    \item represented women's genitalia as underdeveloped versions of men's
    \item argued for male superiority over women
\end{enumerate}
\item In contemporary societies, social institutions are:
\begin{enumerate}[label=(\alph*)]
    \item highly specialized, interrelated sets of social practices
    \item disorganized social relations in a postmodern world
    \item virtual communities in cyberspace
    \item no longer relevant to sociology
\end{enumerate}
\item Durkheim defined social facts as:
\begin{enumerate}[label=(\alph*)]
    \item ways of acting, thinking, and feeling that are collective and social in origin
    \item the way scientists construct knowledge in a social context
    \item data collected about social phenomena that are proven to be correct
    \item ideas and theories that have no basis in the external, physical world
\end{enumerate}
\end{enumerate}

\section*{True / False (10 marks)}
\textit{Answer True or False}
\begin{enumerate}[label=\arabic*.]
\setcounter{enumi}{5}
\item True or False: Disciplinary power in 19 th century factories was largely enforced by the exclusion of women from waged labor roles.
\item True or False: The post-war welfare state of 1948 did not aim to provide a minimum wage for its citizens.
\item True or False: In Marx's theory, the 'mode of production' is defined as the average measure of productivity under capitalism.
\item True or False: Sociology is a social science because its theories are logical, empirical, and subject to scrutiny by peers.
\item True or False: Marx believed that religion would fade away following a socialist revolution that eliminated the necessity for capitalist ideology.
\end{enumerate}

\section*{Long Form Response (20 marks)}
\textit{Answer each question with a clear and complete explanation.}
\begin{enumerate}[label=\arabic*.]
\setcounter{enumi}{10}
\item Discuss Wallerstein's concept of 'world system' in the context of the capitalist world economy, focusing on its implications for the division of labor across diverse ethnic and cultural groups.
\item Discuss the policies pursued by the New Labour government in relation to local education authorities, focusing on the implications of their support for failing institutions.
\end{enumerate}

\vfill
\noindent\textit{End of Paper}
\end{document}